\chapter{TIẾN ĐỘ, ĐÁNH GIÁ VÀ HƯỚNG HOÀN THIỆN}
\label{chap:tiendo}

\section{Tiến độ thực hiện}

Bảng~\ref{tab:tiendo} đối chiếu kế hoạch trong README của nhóm với thời điểm thực hiện, lấy theo lịch sử
commit trên nhánh \code{feat/tv2-init-local-search}.

\begin{table}[H]
\centering
\caption{Tiến độ phần việc TV2}
\label{tab:tiendo}
\begin{tabularx}{\textwidth}{@{}p{3.3cm}Xp{2.2cm}@{}}
\toprule
\textbf{Giai đoạn (kế hoạch)} & \textbf{Công việc} & \textbf{Thực hiện} \\
\midrule
Giai đoạn 2 \newline (12/11--18/11) &
  Khảo sát mã nguồn; cài đặt \code{init\_smart}, \code{local\_search}, lớp dán \code{tv2\_hooks};
  170 test đơn vị; harness NSGA-II tạm vì \code{nsga2.py} lúc đó chưa có; mốc hồi quy khi tắt cờ;
  thí nghiệm sơ bộ trên harness; tài liệu hợp đồng giao diện & 12/11 \\
\midrule
Giai đoạn 3 \newline (19/11--25/11) &
  Quy trình ghép nhánh và xử lý xung đột \code{config.yaml}; ghép thử với nhánh của TV4 trong một nhánh
  tạm và kiểm tra tìm kiếm cục bộ chạy cùng DP split & 12/11 \\
 & Ghép nhánh chính mới nhất (đã có \code{nsga2.py} của TV1, cải tiến của TV3 và TV4) vào nhánh TV2;
  gắn hai cải tiến vào \code{nsga2.py} thật; thêm test tích hợp; 432 test pass & 27/11 \\
\midrule
Giai đoạn 4 \newline (26/11--02/12) &
  Chạy lại toàn bộ thực nghiệm trên thuật toán thật (440 lần chạy); sinh bảng và hình tự động; viết
  báo cáo này & 27/11 \\
\bottomrule
\end{tabularx}
\end{table}

Phần việc của giai đoạn 2 hoàn thành sớm. Phần gắn vào thuật toán thật bị chậm so với kế hoạch: thuật
toán gốc của TV1 được ghép vào nhánh chính ngày 13/11, nhưng nhánh TV2 chỉ cập nhật nhánh chính vào ngày
27/11, khi nhóm trưởng đã ghép xong các nhánh của TV3 và TV4. Trong khoảng thời gian đó, thực nghiệm của
TV2 chạy trên harness tạm; các số liệu trong báo cáo này đều đã được chạy lại trên thuật toán thật. Việc
ghép bốn nhánh vào nhánh chính ở giai đoạn 3 do nhóm trưởng thực hiện; TV2 chuẩn bị quy trình, mốc hồi
quy và kiểm tra tương thích.

\section{Đối chiếu với yêu cầu}

\begin{longtable}{@{}p{8.4cm}p{1.8cm}p{4.3cm}@{}}
\caption{Đối chiếu kết quả với các yêu cầu của phần việc TV2}\label{tab:doichieu}\\
\toprule
\textbf{Yêu cầu} & \textbf{Trạng thái} & \textbf{Bằng chứng} \\
\midrule
\endfirsthead
\toprule
\textbf{Yêu cầu} & \textbf{Trạng thái} & \textbf{Bằng chứng} \\
\midrule
\endhead
Khởi tạo trộn NN, theo khung giờ và ngẫu nhiên theo tỉ lệ trong cấu hình & Đạt & Chương~\ref{chap:init} \\
Hai thứ tự tất định có biến thể, quần thể không trùng lặp & Đạt & mục 3.4, 3.5; test \\
Mọi cá thể là hoán vị hợp lệ; tỉ lệ các nhóm đúng & Đạt & 91 test của \code{init\_smart} \\
2-opt và Or-opt trong từng tuyến sau đột biến & Đạt & Chương~\ref{chap:ls} \\
Định nghĩa rõ thế nào là tốt hơn khi có ba mục tiêu & Đạt & mục 4.4: Pareto hoặc tổng có trọng số \\
Giới hạn số vòng lặp của tìm kiếm cục bộ & Đạt & ba trần cứng, có test \\
Kiểm tra việc cắt lại tuyến sau khi ghép (lệch biên) & Đạt & mục~\ref{sec:drift}; phản ví dụ trong test \\
Mỗi tuyến giữ nguyên tập khách; không bước nào làm tuyến tệ hơn theo tiêu chí & Đạt & test \\
Tham số trong \code{config.yaml}, cờ mặc định tắt & Đạt & mục 3.1, 4.1 \\
Cờ tắt không dùng bộ sinh ngẫu nhiên; ra đúng thuật toán gốc & Đạt & mục~\ref{sec:regress}: trùng từng bit \\
Mọi lần đánh giá thêm được đếm vào \code{n\_evals} & Đạt & công thức~\ref{eq:nevals}, có test \\
Chạy 10 seed trên ít nhất 3 instance, so HV với thuật toán gốc & Đạt & Chương~\ref{chap:thucnghiem}: 4 instance \\
So sánh cùng số lần đánh giá & Đạt & cấu hình gốc chạy dài \\
Chuẩn bị mốc hồi quy và quy trình ghép nhánh & Đạt & \code{tv2\_merge\_playbook.md} \\
Biểu đồ front Pareto 3D, hội tụ HV, hộp HV và IGD & Đạt & hình~\ref{fig:hoitu}--\ref{fig:pareto} \\
Ghép bốn nhánh vào nhánh chính & Nhóm trưởng thực hiện & lịch sử nhánh chính 27/11 \\
Kiểm định thống kê chính thức cho toàn bộ ablation & Chưa & thuộc phần việc TV3 \\
\bottomrule
\end{longtable}

\section{Đánh giá}

\paragraph{Điểm làm được.} Hai cải tiến chạy được trên thuật toán thật, đúng các quy tắc chung của nhóm, có
bộ test bao phủ các trường hợp biên và một bằng chứng mạnh rằng nhánh cờ tắt không đổi kết quả: trùng từng
bit với mốc do TV4 sinh từ mã chưa sửa. Thực nghiệm có tách seed tinh chỉnh khỏi seed báo cáo và có phép so
ở cùng số lần đánh giá, nên kết luận không bị thiên vị cho tìm kiếm cục bộ. Khởi tạo thông minh cho cải
thiện rõ và ổn định trên cả bốn instance.

\paragraph{Điểm còn hạn chế.}
\begin{itemize}
  \item Tìm kiếm cục bộ tốn nhiều lần đánh giá (gấp 3,3 đến 8,1 lần thuật toán gốc) và chiếm 71 đến 89\%
        thời gian chạy. Ở cùng số lần đánh giá, nó chỉ có lợi khi đi cùng khởi tạo thông minh và chỉ trên
        hai trong bốn instance.
  \item Với chính sách mặc định \code{reject}, phần lớn ứng viên bị loại vì lệch biên tuyến, nên tìm kiếm
        cục bộ thuần trong tuyến gần như không có tác dụng với cách cắt \code{capacity\_time}.
  \item Mỗi ứng viên được chấm lại bằng toàn bộ bộ giải mã, $O(n)$, dù chỉ một tuyến thay đổi.
  \item Thực nghiệm mới trên bốn instance, chưa đủ chín instance Solomon của bài báo.
  \item Lợi thế của khởi tạo thông minh giảm khi thuật toán được chạy lâu hơn.
\end{itemize}

\section{Hướng hoàn thiện và cải tiến}

\begin{enumerate}
  \item \textbf{Chọn chính sách lệch biên mặc định.} Số liệu cho thấy \code{evaluate\_actual} tốt hơn
        \code{reject} trên các instance lớn và chưa kém ở đâu có ý nghĩa thống kê. Nên đề xuất với TV4 (người
        phụ trách tinh chỉnh) đổi mặc định sau khi chạy đủ chín instance.
  \item \textbf{Đánh giá tăng dần.} Thay vì giải mã lại cả hoán vị, chỉ tính lại tuyến bị thay đổi và kiểm
        tra chỗ cắt ở hai đầu tuyến. Giờ đến trên các tuyến khác không phụ thuộc tuyến này, vì mỗi xe xuất
        phát độc lập từ kho, nên phần tiết kiệm có thể lớn với $n = 100$.
  \item \textbf{Bước đi giữa các tuyến.} Thêm relocate, swap và 2-opt* giữa hai tuyến. Với cách cắt theo
        thời gian, chuyển khách giữa các tuyến có thể hợp với cấu trúc bài toán hơn bước đi trong tuyến.
  \item \textbf{Chọn cá thể để tìm kiếm cục bộ.} Chỉ áp dụng cho con thuộc front tốt nhất, hoặc giảm $p_{ls}$
        theo thời gian, để giảm số lần đánh giá ở các thế hệ đầu khi quần thể còn kém.
  \item \textbf{Khởi tạo đa dạng hơn.} Dùng láng giềng gần nhất ngẫu nhiên hoá (chọn trong $k$ khách gần
        nhất) hoặc heuristic chèn có xét khung thời gian kiểu Solomon I1 thay cho phép hoán đổi ngẫu nhiên.
  \item \textbf{Thực nghiệm đầy đủ.} Chạy trên chín instance Solomon và case29, so cả theo thời gian chạy,
        và để TV3 áp dụng kiểm định có hiệu chỉnh cho nhiều phép so sánh.
\end{enumerate}
